% =================================================================
%  CSE 815 -- Machine Learning
%  Module 2, Week 6, Part 2
%  The Machine Learning Development Process: Error Analysis, Data,
%  Transfer Learning, Deployment, Fairness, and Skewed Classes
% =================================================================
\documentclass[aspectratio=169,10pt]{beamer}
\usepackage{cse815}

\title[CSE 815 -- Week 6, Part 2]{The Machine Learning Development Process}
\subtitle{Error Analysis, Adding Data, Transfer Learning, Deployment, Fairness, and Precision--Recall}
\author{Rokan Uddin Faruqui\\ \small Professor, Department of CSE}
\institute{\small University of Chittagong \\ \texttt{rufaruqui@cu.ac.bd}}
\date[Week 6 / Part 2]{Week 6 $\cdot$ Part 2}

\begin{document}

\begin{frame}[plain,noframenumbering]
  \titlepage
\end{frame}

% ---------------------------------------------------------------
\begin{frame}{From diagnosis to a development process}
  \begin{columns}[T,onlytextwidth]
    \begin{column}{0.5\textwidth}
      \begin{block}{Part 1 gave us}
        \footnotesize
        \begin{itemize}
          \item Train / cross-validation / test: fit, choose, report.
          \item Bias gap and variance gap, measured against a baseline.
          \item Learning curves: does more data help?
        \end{itemize}
      \end{block}
      \footnotesize
      Bias and variance say \emph{how much} a model is wrong. They do not say \emph{which}
      examples it gets wrong, where more data should come from, or whether the model should
      be deployed at all. Today is about the rest of the process.
    \end{column}
    \begin{column}{0.46\textwidth}
      \begin{keyidea}[Today]
        \footnotesize
        \begin{enumerate}
          \item The iterative loop, and error analysis
          \item Adding data: targeted collection, augmentation, synthesis
          \item Transfer learning
          \item The full cycle of a project; fairness, bias and ethics
          \item Skewed classes: precision, recall and $F_1$
        \end{enumerate}
      \end{keyidea}
    \end{column}
  \end{columns}
\end{frame}

% ===============================================================
\section{The Iterative Loop and Error Analysis}

\begin{frame}{The iterative loop of ML development}
  \begin{columns}[T,onlytextwidth]
    \begin{column}{0.48\textwidth}
      \centering
      \begin{tikzpicture}[
          st/.style={draw=cublue,fill=culight,rounded corners,align=center,font=\scriptsize,
                     minimum width=2.6cm,minimum height=9mm}]
        \node[st] (a) at (90:1.55) {choose architecture\\(model, data, $\lambda$)};
        \node[st,draw=cuamber,fill=cuamber!12] (t) at (-30:1.85) {train model};
        \node[st,draw=cuteal,fill=cuteal!12] (d) at (210:1.85) {diagnostics\\(bias, variance,\\error analysis)};
        \foreach \u/\v in {a/t,t/d,d/a}{
          \draw[-{Stealth[length=2.6mm]},cugray,very thick] (\u) to[bend left=25] (\v);}
      \end{tikzpicture}
    \end{column}
    \begin{column}{0.48\textwidth}
      \footnotesize
      \begin{block}{Example: a spam classifier}
        Features: $x_j=1$ if word $j$ of a 10\,000-word list appears in the email. Ideas
        for improving it:
        \begin{itemize}
          \item collect more data (``honeypot'' accounts that attract spam);
          \item features from the email's routing headers;
          \item features from the body: punctuation, ``w4tches'' for ``watches'';
          \item algorithms to detect deliberate misspellings.
        \end{itemize}
      \end{block}
      \begin{keyidea}
        Every turn of the loop costs time. Good diagnostics make each turn count.
      \end{keyidea}
    \end{column}
  \end{columns}
\end{frame}

\begin{frame}{Error analysis: look at the mistakes}
  \begin{columns}[T,onlytextwidth]
    \begin{column}{0.52\textwidth}
      \footnotesize
      \begin{block}{Illustrative: 500 of 5000 validation emails misclassified}
        Pick 100 of the 500 at random and tag each by hand:\par
        \vspace{1mm}
        \renewcommand{\arraystretch}{1.15}
        \begin{tabular}{@{}lr@{}}
          \toprule
          category & count\\
          \midrule
          pharmaceutical spam & $21$\\
          phishing (stealing passwords) & $18$\\
          unusual email routing & $7$\\
          spam inside an embedded image & $5$\\
          deliberate misspellings & $3$\\
          \bottomrule
        \end{tabular}\par
        \vspace{1mm}
        Categories may overlap: one email can be both pharmaceutical and misspelt.
      \end{block}
    \end{column}
    \begin{column}{0.44\textwidth}
      \footnotesize
      \begin{keyidea}[What the table decides]
        A perfect misspelling detector would fix at most $3$ in $100$ mistakes --- about
        $15$ of the $500$. Pharmaceutical spam and phishing together account for $39$:
        that is where effort should go (more pharmaceutical examples; features for URLs in
        phishing emails).
      \end{keyidea}
      Why 100? Enough to see the large categories clearly; few enough to read in an hour or
      two.
      \begin{pitfall}[Limits]
        \footnotesize Easy for tasks people do well (reading email, recognising objects);
        much harder when people cannot judge either, such as predicting which advert a
        user will click.
      \end{pitfall}
    \end{column}
  \end{columns}
\end{frame}

% ===============================================================
\section{Adding Data}

\begin{frame}{More data --- of the right kind}
  \begin{columns}[T,onlytextwidth]
    \begin{column}{0.5\textwidth}
      \begin{block}{Targeted collection}
        \footnotesize
        ``More of everything'' is slow and expensive. If error analysis says pharmaceutical
        spam is the problem, collect \emph{pharmaceutical spam}: for instance, search the
        unlabelled mail you already have for it and label only that.
      \end{block}
      \begin{block}{Two views of a project}
        \footnotesize
        \textbf{Model-centric}: hold the data fixed, improve the algorithm.\\[2pt]
        \textbf{Data-centric}: hold a good algorithm fixed, improve the data. With strong
        off-the-shelf networks, the second is often the faster route to better
        performance.
      \end{block}
    \end{column}
    \begin{column}{0.46\textwidth}
      \footnotesize
      \begin{defnbox}[Three ways to add data]
        \small
        \begin{enumerate}
          \item \textbf{Collect} more real examples, targeted by error analysis.
          \item \textbf{Augment}: create new examples by modifying existing ones, keeping
                the label.
          \item \textbf{Synthesise}: create new examples from scratch.
        \end{enumerate}
      \end{defnbox}
      \begin{keyidea}
        All three serve a model with \emph{high variance}. Part 1's learning curve tells
        you whether that is the model you have.
      \end{keyidea}
    \end{column}
  \end{columns}
\end{frame}

\begin{frame}{Data augmentation}
  \begin{columns}[T,onlytextwidth]
    \begin{column}{0.52\textwidth}
      \centering {\scriptsize One labelled ``A'' becomes many}\par
      \vspace{1mm}
      \begin{tikzpicture}[letter/.style={font=\sffamily\bfseries\Huge,cublue}]
        \node[draw=cuamber,thick,minimum size=1.3cm,letter] at (0,0) {A};
        \node[font=\tiny,cuamber] at (0,-0.9) {original};
        \foreach \x/\lab in {1.8/rotated,3.4/stretched,5.0/shrunk,6.6/sheared}{
          \node[draw=cugray!50,minimum size=1.3cm] at (\x,0) {};
          \node[font=\tiny,cugray] at (\x,-0.9) {\lab};}
        \node[letter,rotate=15] at (1.8,0) {A};
        \node[letter,xscale=1.4] at (3.4,0) {A};
        \node[letter,scale=0.65] at (5.0,0) {A};
        \node[letter,xslant=0.4] at (6.6,0) {A};
        \node[draw=cugray!50,minimum size=1.3cm] at (3.4,-2.1) {};
        \node[letter,cublue!45] at (3.4,-2.1) {A};
        \pgfmathsetseed{6}
        \foreach \i in {1,...,40}{
          \pgfmathsetmacro{\px}{3.4+1.1*(rnd-0.5)}
          \pgfmathsetmacro{\py}{-2.1+1.1*(rnd-0.5)}
          \fill[cugray] (\px,\py) circle (0.6pt);}
        \node[font=\tiny,curose] at (3.4,-3.0) {random pixel noise: usually useless};
      \end{tikzpicture}
    \end{column}
    \begin{column}{0.44\textwidth}
      \footnotesize
      \begin{block}{Other domains}
        \textbf{Speech}: add background noise --- a crowd, a car, a bad phone line --- to
        a clean recording. \textbf{Text}: replace words with synonyms. Each modified copy
        keeps the original label.
      \end{block}
      \begin{keyidea}[The rule]
        A distortion helps only if it represents variation you expect in the
        \emph{test} set. Rotated letters and noisy calls occur in practice; per-pixel
        random noise does not, and adds little.
      \end{keyidea}
      \begin{pitfall}[Do not change the label]
        \footnotesize Rotate a ``6'' by $180^\circ$ and it is a ``9''. Augmentations must
        preserve what the label means.
      \end{pitfall}
    \end{column}
  \end{columns}
\end{frame}

\begin{frame}{Data synthesis}
  \begin{columns}[T,onlytextwidth]
    \begin{column}{0.5\textwidth}
      \begin{block}{Photo OCR}
        \footnotesize
        Reading text in photographs of shop signs and number plates. Real labelled crops are
        scarce. Synthesise them: render characters in many fonts, sizes, colours and
        contrasts, on varied backgrounds, and crop.
      \end{block}
      \centering
      \begin{tikzpicture}
        \foreach \x/\f/\c in {0/\rmfamily/cublue,1.0/\sffamily/curose,2.0/\ttfamily/cuteal,
                              3.0/\itshape/cuamber,4.0/\bfseries/cugray}{
          \node[draw=cugray!40,fill=culight,minimum size=8mm,font=\f\LARGE,text=\c]
            at (\x,0) {A};
        }
      \end{tikzpicture}\par
      {\scriptsize synthetic training images of the letter ``A''}
    \end{column}
    \begin{column}{0.46\textwidth}
      \footnotesize
      \begin{block}{When synthesis works}
        When you can write a program whose output looks like real inputs, with the label
        known by construction: rendered text, simulated sensor readings, game-engine images
        for self-driving.
      \end{block}
      \begin{pitfall}[The gap between synthetic and real]
        \footnotesize A model can learn the renderer's quirks instead of the task. Always
        validate on \emph{real} data, and check with error analysis that the synthetic data
        fixes the categories it was made for.
      \end{pitfall}
    \end{column}
  \end{columns}
\end{frame}

% ===============================================================
\section{Transfer Learning}

\begin{frame}{Using data from a different task}
  \begin{columns}[T,onlytextwidth]
    \begin{column}{0.54\textwidth}
      \centering
      \begin{tikzpicture}[scale=1.35,transform shape,node distance=2.5mm,
          lay/.style={draw=cublue,fill=culight,rounded corners,font=\tiny,
                      minimum height=5mm,inner sep=2pt,align=center}]
        \node[font=\scriptsize\bfseries,cublue] at (2.4,0.75) {1. Supervised pre-training};
        \node[lay,draw=cuteal,fill=cuteal!12] (x) at (0,0) {$\vx$: 784};
        \node[lay,right=of x] (h1) {128};
        \node[lay,right=of h1] (h2) {64};
        \node[lay,right=of h2,draw=cuamber,fill=cuamber!15] (o) {1000 classes};
        \foreach \a/\b in {x/h1,h1/h2,h2/o}{\draw[synapse] (\a)--(\b);}
        \node[font=\tiny,cugray] at (2.4,-0.55) {1 million labelled images: cats, cars, people, \dots};
        \node[font=\scriptsize\bfseries,cublue] at (2.4,-1.25) {2. Fine-tuning};
        \node[lay,draw=cuteal,fill=cuteal!12] (x2) at (0,-2.0) {$\vx$: 784};
        \node[lay,right=of x2] (g1) {128};
        \node[lay,right=of g1] (g2) {64};
        \node[lay,right=of g2,draw=curose,fill=curose!10] (o2) {10 digits};
        \foreach \a/\b in {x2/g1,g1/g2,g2/o2}{\draw[synapse] (\a)--(\b);}
        \node[font=\tiny,cugray] at (2.4,-2.55) {small labelled set of digits};
        \draw[cuteal,thick,dashed,-{Stealth}] (h1.south) -- node[left,font=\tiny]{copy} (g1.north);
        \draw[cuteal,thick,dashed,-{Stealth}] (h2.south) -- (g2.north);
      \end{tikzpicture}
    \end{column}
    \begin{column}{0.42\textwidth}
      \footnotesize
      \begin{block}{Two options for step 2}
        \textbf{Option 1}: train only the new output layer; keep the copied layers fixed.
        Best when the new data set is very small.\\[2pt]
        \textbf{Option 2}: train every layer, starting from the copied values. Better when
        you have more data.
      \end{block}
      \begin{block}{Parameters}
        Copied: $100\,480+8\,256=108\,736$. New output layer: $64\cdot10+10=650$ (it
        replaces $65\,000$). Option 1 trains $650$ of $109\,386$.
      \end{block}
    \end{column}
  \end{columns}
\end{frame}

\begin{frame}{Why transfer learning works}
  \begin{columns}[T,onlytextwidth]
    \begin{column}{0.5\textwidth}
      \begin{keyidea}[Generic early features]
        Week 4: the early layers of an image network detect edges, then corners and curves,
        then parts. These are useful for almost \emph{any} image task, so layers learned on
        a million photographs give a head start on a few hundred X-rays or digits.
      \end{keyidea}
      \begin{pitfall}[Same type of input]
        \footnotesize Pre-train on images for an image task, on audio for a speech task, on
        text for a language task. A network pre-trained on photographs offers nothing to a
        spreadsheet of loan applications.
      \end{pitfall}
    \end{column}
    \begin{column}{0.46\textwidth}
      \footnotesize
      \begin{block}{In practice}
        \begin{enumerate}
          \item Download a network pre-trained on a large data set with the same input type
                --- often published by researchers or companies.
          \item Replace the output layer to suit your classes.
          \item Fine-tune on your own data (option 1 or 2).
        \end{enumerate}
      \end{block}
      \begin{block}{The same idea, at scale}
        Large language models such as GPT and BERT are pre-trained on vast amounts of text,
        then fine-tuned for a specific task --- including tasks in Bangla.
      \end{block}
    \end{column}
  \end{columns}
\end{frame}

% ===============================================================
\section{The Full Cycle, and Fairness}

\begin{frame}{The full cycle of a machine learning project}
  \centering
  \begin{tikzpicture}[
      st/.style={draw=cublue,fill=culight,rounded corners,align=center,font=\scriptsize,
                 minimum width=2.6cm,minimum height=11mm}]
    \node[st] (s) at (0,0) {\textbf{Scope}\\define the project};
    \node[st] (c) at (3.4,0) {\textbf{Collect data}\\define and collect};
    \node[st,draw=cuamber,fill=cuamber!12] (t) at (6.8,0) {\textbf{Train model}\\train, error analysis,\\iterate};
    \node[st,draw=cuteal,fill=cuteal!12] (d) at (10.2,0) {\textbf{Deploy}\\deploy, monitor,\\maintain};
    \draw[synapse,very thick] (s) -- (c);
    \draw[synapse,very thick] (c) -- (t);
    \draw[synapse,very thick] (t) -- (d);
    \draw[synapse,thick,curose] (t.south) to[bend left=35] node[below,font=\tiny,fill=white,inner sep=1pt]{need more data} (c.south);
    \draw[synapse,thick,curose] (d.south) to[bend left=30] node[below,font=\tiny,fill=white,inner sep=1pt]{performance drifts: retrain} (t.south);
  \end{tikzpicture}\par
  \vspace{2mm}
  \begin{columns}[T,onlytextwidth]
    \begin{column}{0.5\textwidth}
      \begin{block}{Deployment}
        \footnotesize
        The model runs on an \textbf{inference server}. A mobile app sends an input --- a
        voice query, a transaction --- through an API call and receives a prediction.
        Engineering concerns then dominate: reliable prediction at scale, logging inputs
        and outputs, monitoring, and updating the model.
      \end{block}
    \end{column}
    \begin{column}{0.46\textwidth}
      \begin{keyidea}[MLOps]
        \footnotesize Machine-learning operations: the practice of building, deploying and
        maintaining ML systems. Logged data from deployment is often the best new training
        data --- and the first sign that the world has changed (\textbf{data drift}).
      \end{keyidea}
    \end{column}
  \end{columns}
\end{frame}

\begin{frame}{Fairness, bias and ethics}
  \begin{columns}[T,onlytextwidth]
    \begin{column}{0.5\textwidth}
      \footnotesize
      \begin{block}{When systems have done harm}
        \begin{itemize}
          \item A hiring tool that learned to rank women's applications lower, from past
                hiring decisions.
          \item Face recognition that is markedly less accurate for darker-skinned faces.
          \item Loan approval that disadvantages groups under-represented in the training
                data.
        \end{itemize}
      \end{block}
      \begin{block}{Adverse uses}
        Deepfakes made without consent; recommender systems that spread toxic content
        because it is engaging; fraud and fake reviews generated at scale.
      \end{block}
    \end{column}
    \begin{column}{0.46\textwidth}
      \footnotesize
      \begin{keyidea}[Bias in, bias out]
        ``Bias'' here is not the bias of Part 1. A model trained on historical decisions
        learns the patterns in those decisions --- including unfair ones --- and applies
        them at scale.
      \end{keyidea}
      \begin{block}{Practical guidelines}
        \begin{enumerate}
          \item A diverse team brainstorms what could go wrong, especially for
                vulnerable groups.
          \item Search the literature for standards in your field.
          \item Audit for harm \emph{before} deployment, e.g.\ accuracy per group.
          \item Plan mitigation (e.g.\ roll back) and keep monitoring.
        \end{enumerate}
      \end{block}
    \end{column}
  \end{columns}
\end{frame}

% ===============================================================
\section{Skewed Classes: Precision and Recall}

\begin{frame}{When accuracy misleads}
  \begin{columns}[T,onlytextwidth]
    \begin{column}{0.5\textwidth}
      \footnotesize
      \begin{block}{Screening for a rare disease}
        $1000$ patients in the validation set, $25$ of them ill ($y=1$).
        \begin{itemize}
          \item A trained classifier: accuracy $98.5\%$.
          \item \texttt{print("y = 0")} for everyone: accuracy $97.5\%$.
        \end{itemize}
        The program that never finds a single patient is nearly as ``accurate''.
      \end{block}
      \begin{block}{Confusion matrix of the classifier}
        \centering
        \renewcommand{\arraystretch}{1.25}
        \begin{tabular}{@{}l|cc@{}}
          & actual $1$ & actual $0$\\
          \hline
          predicted $1$ & TP $=15$ & FP $=5$\\
          predicted $0$ & FN $=10$ & TN $=970$\\
        \end{tabular}
      \end{block}
    \end{column}
    \begin{column}{0.46\textwidth}
      \footnotesize
      \begin{defnbox}[Precision and recall]
        \small
        \[ \text{precision}=\frac{\text{TP}}{\text{TP}+\text{FP}}=\frac{15}{20}=0.75 \]
        of the patients we flagged, the fraction actually ill.
        \[ \text{recall}=\frac{\text{TP}}{\text{TP}+\text{FN}}=\frac{15}{25}=0.60 \]
        of the ill patients, the fraction we found.
      \end{defnbox}
      \begin{keyidea}
        ``Always $y=0$'' has recall $0$ and an undefined precision ($0/0$): the two numbers
        expose it immediately. Convention: $y=1$ is the rare class.
      \end{keyidea}
    \end{column}
  \end{columns}
\end{frame}

\begin{frame}{Trading precision for recall with the threshold}
  \begin{columns}[T,onlytextwidth]
    \begin{column}{0.48\textwidth}
      \centering {\scriptsize Precision--recall curve as the threshold varies}\par
      \begin{tikzpicture}
        \begin{axis}[width=6.4cm,height=4.5cm,xmin=0,xmax=1.02,ymin=0,ymax=1.02,
          xlabel={recall},ylabel={precision},tick label style={font=\scriptsize},
          label style={font=\scriptsize}]
          \addplot[cublue,very thick,smooth] coordinates
            {(0.999,0.117)(0.997,0.125)(0.994,0.138)(0.988,0.155)(0.977,0.178)(0.960,0.210)
             (0.933,0.252)(0.894,0.305)(0.841,0.371)(0.773,0.449)(0.691,0.535)(0.599,0.624)
             (0.500,0.709)(0.401,0.785)(0.309,0.847)(0.227,0.894)(0.159,0.929)(0.106,0.953)
             (0.067,0.970)(0.040,0.981)(0.023,0.988)};
          \addplot[only marks,mark=*,mark size=2pt,curose] coordinates
            {(0.841,0.371)(0.627,0.597)(0.461,0.741)(0.301,0.851)(0.116,0.949)};
          \node[font=\tiny,curose,anchor=south west] at (axis cs:0.841,0.371) {0.1};
          \node[font=\tiny,curose,anchor=south west] at (axis cs:0.627,0.597) {0.3};
          \node[font=\tiny,curose,anchor=south west] at (axis cs:0.461,0.741) {0.5};
          \node[font=\tiny,curose,anchor=south west] at (axis cs:0.301,0.851) {0.7};
          \node[font=\tiny,curose,anchor=south west] at (axis cs:0.116,0.949) {0.9};
        \end{axis}
      \end{tikzpicture}\par
      {\scriptsize Illustrative: 10\% positives, scores $\mathcal{N}(2,1)$ vs.\
      $\mathcal{N}(0,1)$; labels are thresholds on $P(y{=}1\mid\vx)$.}
    \end{column}
    \begin{column}{0.48\textwidth}
      \footnotesize
      \begin{block}{Predict $\yhat=1$ if $f(\vx)\ge$ threshold}
        \renewcommand{\arraystretch}{1.0}
        \begin{tabular}{@{}cccc@{}}
          \toprule
          threshold & precision & recall & $F_1$\\
          \midrule
          $0.1$ & $0.371$ & $0.841$ & $0.515$\\
          $0.3$ & $0.597$ & $0.627$ & $\mathbf{0.612}$\\
          $0.5$ & $0.741$ & $0.461$ & $0.568$\\
          $0.7$ & $0.851$ & $0.301$ & $0.444$\\
          $0.9$ & $0.949$ & $0.116$ & $0.206$\\
          \bottomrule
        \end{tabular}
      \end{block}
      \textbf{Raise} it: predict $1$ only when confident --- higher precision, lower recall
      (before an invasive treatment). \textbf{Lower} it: miss fewer cases --- higher recall,
      lower precision (when a miss is dangerous and a follow-up test is cheap).
      \begin{pitfall}[A choice, not a fit]
        \footnotesize Set the threshold from the costs of the two kinds of error.
      \end{pitfall}
    \end{column}
  \end{columns}
\end{frame}

\begin{frame}{Combining precision and recall: the $F_1$ score}
  \begin{columns}[T,onlytextwidth]
    \begin{column}{0.5\textwidth}
      \footnotesize
      \begin{block}{Three algorithms --- which is best?}
        \renewcommand{\arraystretch}{1.15}
        \begin{tabular}{@{}lcccc@{}}
          \toprule
          & $P$ & $R$ & average & $F_1$\\
          \midrule
          algorithm 1 & $0.5$  & $0.4$ & $0.45$ & $\mathbf{0.444}$\\
          algorithm 2 & $0.7$  & $0.1$ & $0.40$ & $0.175$\\
          algorithm 3 & $0.02$ & $1.0$ & $\mathbf{0.51}$ & $0.0392$\\
          \bottomrule
        \end{tabular}
      \end{block}
      Algorithm 3 predicts $y=1$ for almost everyone: perfect recall, useless precision.
      The plain average ranks it first.
    \end{column}
    \begin{column}{0.46\textwidth}
      \begin{defnbox}[$F_1$ score]
        \small
        \[ F_1=\frac{1}{\tfrac12\bigl(\tfrac1P+\tfrac1R\bigr)}=\frac{2PR}{P+R} \]
        The \textbf{harmonic mean} of $P$ and $R$. It is dominated by the smaller of the
        two, so a classifier must do well on both to score well.
      \end{defnbox}
      \footnotesize
      \textbf{Check:} algorithm 1, $2(0.5)(0.4)/0.9=0.444$; the screening classifier,
      $2(0.75)(0.6)/1.35=0.667$.\par
      \vspace{1mm}
      In the threshold table, $F_1$ peaks at about $0.31$ ($F_1=0.612$) --- not at the
      default $0.5$.
    \end{column}
  \end{columns}
\end{frame}

% ===============================================================
\section{Wrap-up and Assessment}

\begin{frame}{Summary}
  \begin{columns}[T,onlytextwidth]
    \begin{column}{0.5\textwidth}
      \begin{block}{What to remember}
        \footnotesize
        \begin{itemize}
          \item Development is a loop: choose, train, diagnose. Diagnostics decide the next
                turn.
          \item Error analysis: tag a sample of the mistakes by hand; fix the largest
                categories first.
          \item Add data where error analysis points; augment with realistic distortions;
                synthesise when you can render the input.
          \item Transfer learning: pre-train on a large related task, replace the output
                layer, fine-tune.
          \item Deploy, monitor and retrain; audit for harm before and after.
          \item Skewed classes: use precision, recall and $F_1$, not accuracy.
        \end{itemize}
      \end{block}
    \end{column}
    \begin{column}{0.46\textwidth}
      \begin{block}{Formulas}
        \[ P=\frac{\text{TP}}{\text{TP}+\text{FP}},\qquad R=\frac{\text{TP}}{\text{TP}+\text{FN}} \]
        \[ F_1=\frac{2PR}{P+R} \]
      \end{block}
      \begin{block}{Next}
        \footnotesize Module 3, Week 7: unsupervised learning --- clustering with $k$-means,
        and anomaly detection.
      \end{block}
    \end{column}
  \end{columns}
\end{frame}

\begin{frame}{Descriptive questions}
  \footnotesize
  \begin{enumerate}
    \item Describe the iterative loop of ML development and explain the role of
          diagnostics in it. \hfill\clo{4}
    \item Explain how to carry out an error analysis. Using the spam table, explain why a
          misspelling detector is a poor use of effort. \hfill\clo{6}
    \item Distinguish collecting, augmenting and synthesising data, with one example of
          each. When will none of them help? \hfill\clo{6}
    \item State the rule for choosing data augmentations, and give one augmentation that
          would harm a digit classifier. \hfill\clo{5}
    \item Explain transfer learning, its two fine-tuning options, and when to use each. Why
          must the pre-training input be of the same type? \hfill\clo{6}
    \item Describe the four stages of the full cycle of an ML project, and two reasons a
          deployed model may need retraining. \hfill\clo{5}
  \end{enumerate}
\end{frame}

\begin{frame}{Descriptive questions (continued)}
  \footnotesize
  \begin{enumerate}\setcounter{enumi}{6}
    \item Give two examples of ML systems that caused harm, and four guidelines for reducing
          the risk. \hfill\clo{5}
    \item Explain why accuracy is misleading for skewed classes, using a classifier that
          always predicts $y=0$. \hfill\clo{4}
    \item A fraud model on $10\,000$ transactions gives TP $=40$, FP $=60$, FN $=20$.
          Compute precision, recall, $F_1$ and accuracy. \hfill\clo{6}
    \item Explain how raising the threshold of a logistic classifier changes precision and
          recall, and sketch the precision--recall curve. \hfill\clo{6}
    \item Why is $F_1$ preferred to the arithmetic mean of $P$ and $R$? Compute both for
          $P=0.9$, $R=0.2$. \hfill\clo{5}
    \item For transfer learning from a $784\to128\to64\to1000$ network to a 5-class task,
          how many parameters does option 1 train? \hfill\clo{4}
  \end{enumerate}
\end{frame}

\begin{frame}{Design questions}
  \footnotesize
  \begin{enumerate}
    \item \textbf{Design an error analysis.} A Bangla handwritten-address reader for the
          postal service misreads $8\%$ of envelopes. Design the error analysis: sample
          size, categories you expect, and how each finding would change what you do next.
          \hfill\clo{12}
    \item \textbf{Design an augmentation pipeline.} For a rice-leaf disease classifier
          trained on photos from a few fields, design augmentations that reflect the
          conditions farmers' phones will meet. Name one tempting augmentation you would
          reject, and why. \hfill\clo{10}
    \item \textbf{Design a threshold.} A dengue screening model flags patients for a blood
          test costing Tk~$500$; a missed case costs about Tk~$50\,000$ in later treatment.
          Design how to choose the threshold from validation data, and state what you
          report to the hospital. \hfill\clo{12}
    \item \textbf{Design a fairness audit.} A bank will use a model to pre-screen
          micro-loan applications from rural and urban applicants. Design the audit you
          would run before deployment, the metric per group, and the condition that would
          stop the launch. \hfill\clo{10}
  \end{enumerate}
\end{frame}

\begin{frame}{Reading}
  \begin{block}{Primary text}
    G.~James, D.~Witten, T.~Hastie and R.~Tibshirani,
    \emph{An Introduction to Statistical Learning}, 2nd ed., Springer, 2021.\\[2pt]
    \hfill{\small$\rightarrow$ \textbf{\S4.4.2} --- confusion matrices, changing the
    threshold, and the ROC curve; \textbf{\S10.3.4--10.3.5} --- data augmentation and
    using pre-trained networks.}
  \end{block}
  \begin{block}{Supplementary}
    C.~M.~Bishop, \emph{Pattern Recognition and Machine Learning}, Springer, 2006.\\[2pt]
    \hfill{\small$\rightarrow$ \textbf{\S1.5} --- decision theory: unequal error costs and
    the decision threshold; \textbf{\S5.5.3} --- invariances learned from transformed
    training data.}
  \end{block}
  \vspace{1mm}
  \begin{center}
    \small\color{cugray}
    Lab (CSE 816, Week 6): model evaluation and selection; diagnosing bias and variance.
  \end{center}
\end{frame}

\end{document}
